\section{Capítulo III: Situación Actual, Problemática y Transferencia.}

\subsection{3.1. El problema de la subutilización estructural.}
\subsubsection{Inflexibilidad tarifaria frente a la dinámica del mercado privado.}

Según el informe N° 198-GPRC/2017 del Organismo Supervisor de Inversión Privada en Telecomunicaciones 
(OSIPTEL), la fijación de la tarifa de \$27 dólares por Mbps tuvo origen en el planteamiento 
incial del Ministerio de Transportes y Comunicaciones(MTC) bajo el marco de Ley N° 29904.  
Dicha tarifa regulada debía ser pagada directamente a la empresa concesionaria Azteca Comunicaciones Perú
S.A.C. por parte de los operadores que contrataran la capacidad de transporte.
El MTC partió de supuestos de costos y demanda sobredimensionados, lo que colocó a la red dorsal 
en una clara desventaja competitiva.
Mientras la tarifa estatal estaba congelada por contrato, los operadores privados desplegaron su propia infraestructura y bajaron los precios agresivamente en el mercado, 
llegando a ofrecer costos de un dígito por Mbps. La RDNFO no contempla descuentos por volumen ni flexibilidad geográfica, las empresas privadas solo la utilizan
como ultimo recurso cuando desean extender sus servicios a rutas donde no les saldría a cuenta tender sus propias redes de fibra óptica.

\subsubsection{El quiebre del modelo de Operador Neutro.}

A pesar del gasto global de más de dos mil millones de dólares, a lo largo del tiempo, en 
este macro despliegue técnico, su modelo financiero como ``operador de operadores`` neutro fracasó. Al no alcanzar la demanda proyectada por el MTC, el porcentaje operativo real en 2023 según infobae se estancó en 8\%  de su capacidad total. Además, la brecha fiscal 
conforme pasa el tiempo se vuelve cada vez más inmanejablen, mientras que el mantenimiento
y funcionamiento de la red le cuesta al estado cerca de setenta millones de dólares, la recaudación
efectiva debido al escaso uso llega a los doce millones de dólares.
Como señala el análisis sectorial de Telesemana, bajo este modelo la ganancia es extremadamente 
baja y los beneficiarios reales son mínimos, 
convirtiendo un activo de alta confiabilidad en un ``elefante blanco`` altamente deficitario.
